% !TeX program = xelatex
\documentclass[11pt,UTF8,fontset=fandol]{ctexart}

\usepackage[margin=0.82in]{geometry}
\usepackage{amsmath,amssymb,mathtools}
\usepackage{booktabs,tabularx,array}
\usepackage{float}
\usepackage{enumitem}
\usepackage[dvipsnames]{xcolor}
\usepackage{xurl}
\usepackage{hyperref}

\hypersetup{
  colorlinks=true,
  linkcolor=MidnightBlue,
  citecolor=ForestGreen,
  urlcolor=BrickRed,
  pdftitle={从后续纠正学会第一次回答}
}

\newcolumntype{Y}{>{\raggedright\arraybackslash}X}
\newcolumntype{L}[1]{>{\raggedright\arraybackslash}p{#1}}
\renewcommand{\arraystretch}{1.15}
\setlength{\tabcolsep}{5pt}
\setlist[itemize]{leftmargin=1.35em,itemsep=2pt,topsep=3pt}
\setlist[enumerate]{leftmargin=1.55em,itemsep=3pt,topsep=3pt}

\newcommand{\Dhist}{D_{\mathrm{hist}}}
\newcommand{\Dtrain}{D_{\mathrm{train}}}
\newcommand{\Dholdout}{D_{\mathrm{holdout}}}
\newcommand{\Xholdout}{X_{\mathrm{holdout}}}
\newcommand{\Xsim}{\widetilde X}
\newcommand{\FA}{\mathrm{FA}_{\phi}}
\newcommand{\policy}{\pi_{\theta}}
\newcommand{\basepolicy}{\pi_{\mathrm{base}}}
\newcommand{\lora}{\Delta\theta_u}
\newcommand{\Judge}{\operatorname{Judge}}
\newcommand{\clip}{\operatorname{clip}}
\newcommand{\sg}{\operatorname{stopgrad}}

\begin{document}

\begin{center}
  {\LARGE\bfseries 从后续纠正学会第一次回答\par}
  \vspace{5pt}
  {\large 面向单个用户的离线到在线提示蒸馏与 LoRA 更新\par}
  \vspace{4pt}
  {\small 中文算法设计笔记 \quad 2026 年 8 月 29 日}
\end{center}

\section{先说结论}

单纯的``提示词版本''已经有人做过。DSO 从日志反馈学习提示词策略，
AGP 优化个性化推荐提示词，GEPA 根据执行轨迹改写提示词
\cite{kiyohara2025prompt,wang2025agp,agrawal2025gepa}。因此，不能把
``冻结模型，只优化提示词''写成新的研究点。

现有的 18 动作方案也不适合当前数据。它先把用户偏好压成
\[
  a=(m,\tau,s),\qquad |\mathcal A|=18,
\]
再由手写的 \(\operatorname{Compile}(a,x)\) 生成指令。真实数据却包含自由文本回答、
JSON 工具调用、工具返回和多步状态。18 个标签既不能表示这些动作，也不能从数据中学出
新的行为。

更直接的做法是：
\begin{quote}
用户在后续一轮指出问题后，把这条纠正配回产生错误之前的状态。教师模型在``原状态
+纠正''下给出方向；用户 LoRA 在不看纠正的原状态上学习这个方向。部署后，模型在第一次
回答时就应使用该偏好。
\end{quote}

本文把这一步简称为\textbf{纠正前移}。提示词仍然有用，但只在训练时给教师提供纠正。
部署对象是每位用户自己的 LoRA，不是一份不断变长的提示词。

\section{真实数据告诉了我们什么}

\subsection{当前可用数据}

\begin{table}[H]
  \centering
  \caption{\texttt{/fsx/xzhichao/workspace/oocso/data} 中与本设计直接相关的文件。}
  \begin{tabularx}{\textwidth}{@{}L{0.29\textwidth}L{0.18\textwidth}Y@{}}
    \toprule
    文件 & 规模 & 用途 \\
    \midrule
    \path{tasks_train.jsonl} &
    120 个任务 &
    工具调用训练任务；当前 \path{run.py} 取编号 0--79 做离线训练 \\
    \path{tasks_dev.jsonl} &
    30 个任务 &
    开发阶段检查任务成功率 \\
    \path{tasks_eval.jsonl} &
    30 个任务 &
    最终保留任务，即 \(\Xholdout\) \\
    \path{dump_samples.json} &
    3 个离线会话、2 个在线会话 &
    展示真实字段和当前失败模式，不足以支持最终结论 \\
    \path{preview_pipeline.json} &
    4 个任务、2 条离线线程、3 个在线记录 &
    小规模流程预览 \\
    \path{preview_pi0_student.pt} &
    288 个 LoRA 张量 &
    已训练的用户适配器快照 \\
    \bottomrule
  \end{tabularx}
\end{table}

工具任务分为七类：
\[
  \begin{aligned}
  C=\{&
  \texttt{single\_tool},\texttt{two\_step\_chain},
  \texttt{chain3},\texttt{cross\_rate},\\
  &\texttt{cross\_rate3},\texttt{sum\_two},
  \texttt{error\_correction}\}.
  \end{aligned}
\]
每条任务给出允许使用的工具、答案、误差范围和最大步数。它们可以检查任务是否真的完成，
但本身不包含用户偏好。偏好信号来自同一用户的对话和纠正，不能从工具任务中凭空补写。

\subsection{已经观察到的结构性问题}

\path{dump_samples.json} 中的三条离线会话具有相同结构：
\begin{enumerate}
  \item 第一次回答没有看到偏好提示，回答带有加粗、编号或小标题，奖励为 \(-1\)；
  \item 用户要求去掉这些格式；
  \item 第二次回答看到明确纠正后改成纯文本，奖励为 \(+1\)。
\end{enumerate}

因此，所有正样本都带有明确纠正。普通 GRPO 只能学到``用户提醒后再改''。已有
\(\pi_0\) 在新任务的第一次回答上仍然违规，这与数据结构一致。问题不在于训练轮数太少，
而在于正样本处于错误的状态。

数据还暴露了两个实现问题：有的用户反应没有结束思考块，因而可见文本为空；有的 hint
只剩一个字``和''。这两种记录不能进入纠正前移。

\section{模型、输入和输出}

\subsection{训练什么}

冻结主干为
\[
  \FA=\text{Qwen3-4B-Thinking-2507}.
\]
每位用户 \(u\) 有一个 LoRA 适配器 \(\lora\)。现有检查点对应
\[
  r=16,\qquad \alpha=32,
\]
并插入每层注意力的
\(\texttt{q\_proj},\texttt{k\_proj},\texttt{v\_proj},\texttt{o\_proj}\)。
检查点包含 288 个矩阵和 \(11{,}796{,}480\) 个可训练数。主干参数、工具和环境都不更新。

这里不再额外训练一个 18 类分类器。用户身份由加载哪个 LoRA 决定，所以单用户实验也不需要
用户 ID embedding。

\subsection{输入是什么}

一次决策前的模型输入记为 \(x\)。它是代码中已经存在的
\texttt{state\_messages}/\texttt{state}：
\[
  x=(\text{system},\text{用户任务},\text{此前回答},
  \text{此前工具返回}).
\]
输入只包含本次决策前已经发生的内容。评测第一次回答时，\(x\) 不含当前任务的用户纠正，
也不含人工写入的偏好提示。

\subsection{输出是什么}

输出 \(a=(a_1,\ldots,a_L)\) 是模型实际生成的 token 序列：
\begin{itemize}
  \item 对普通对话，\(a\) 是用户看到的回答；
  \item 对工具任务，\(a\) 是一个
  \texttt{<tool\_call>\{...\}</tool\_call>} JSON 调用；
  \item 多步任务每一步重新形成状态 \(x_t\)，再生成下一条 \(a_t\)。
\end{itemize}

策略为
\begin{equation}
  \policy(a\mid x)
  =\prod_{j=1}^{L}
    \pi_{\phi,\lora}(a_j\mid x,a_{<j}).
  \label{eq:policy}
\end{equation}
动作空间 \(\mathcal A(x)\) 是当前协议允许的 token 序列，不是固定的 18 个标签。

\section{与原 PDF 的符号对应}

\begin{table}[H]
  \centering
  \caption{保留原符号，但改变错误的动作定义。}
  \begin{tabularx}{\textwidth}{@{}L{0.20\textwidth}L{0.33\textwidth}Y@{}}
    \toprule
    符号 & 本文含义 & 与原方案的关系 \\
    \midrule
    \(\Dhist\) &
    同一用户按时间排列的完整会话 &
    保留 \\
    \(\Dtrain,\Dholdout\) &
    过去会话与未来保留会话 &
    保留；必须按完整会话和时间切分 \\
    \(\Xholdout\) &
    最终工具任务和未来用户任务 &
    保留；不参与训练、hint 提取或选参 \\
    \(\Xsim\) &
    可选模拟任务 &
    首个实验设为 \(\varnothing\)，先使用真实任务 \\
    \(C,c\) &
    七类工具任务及之后明确加入的用户任务类 &
    保留；类别来自数据，不由 Judge 临时发明 \\
    \(P_u\) &
    可选的用户偏好摘要 &
    只给教师、Judge 和提示词基线；适配器评测不输入它 \\
    \(a,\mathcal A(x)\) &
    实际回答或工具调用 token 序列 &
    替换 \((m,\tau,s)\) 和 18 动作 \\
    \(y_a\) &
    执行 \(a\) 后的工具返回、最终答案或用户反应 &
    保留为可观察结果 \\
    \(r_a\) &
    由任务结果和真实偏好证据得到的分数 &
    保留，但不由手写动作表产生 \\
    \(A_a=r_a-\bar r\) &
    同一状态候选动作的相对优势 &
    保留 \\
    \(\policy,\basepolicy\) &
    带用户 LoRA 的策略与不带 LoRA 的主干策略 &
    保留 \\
    \(\FA\) &
    冻结主干、工具协议和执行环境 &
    保留 \\
    \bottomrule
  \end{tabularx}
\end{table}

\(\operatorname{Compile}(a,x)\) 被删除。模型输出本身就是可执行动作，不再把标签翻译成
手写提示词。

\section{从历史中构造纠正记录}

对一段历史，保留五项：
\begin{equation}
  e_i=(x_i,a_i^{-},f_i,a_i^{+},o_i).
  \label{eq:record}
\end{equation}
其中 \(x_i\) 是第一次错误回答前的状态，\(a_i^{-}\) 是该回答，\(f_i\) 是用户随后给出的
纠正，\(a_i^{+}\) 是纠正后的回答，\(o_i\) 是任务结果或后续用户反应。

从 \(f_i\) 提取一句具体纠正：
\begin{equation}
  h_i=\operatorname{ExtractCorrection}
      (a_i^{-},f_i,a_i^{+}).
\end{equation}
提取器只能整理用户已经说出的内容，不能补写新的偏好。以下记录丢弃：
\begin{itemize}
  \item \(f_i\) 的可见文本为空；
  \item \(h_i\) 过短、无法执行或只包含泛泛评价；
  \item \(a_i^{+}\) 仍未被接受；
  \item 在有确定性检查器的任务上，纠正后的结果仍然错误。
\end{itemize}

若 \(a_i^{+}\) 脱离上一轮后仍能独立回答原任务，就直接把它作为参考序列
\(\widehat a_i\)。否则，让冻结教师在 \(x_i\oplus h_i\) 下重新生成
\(\widehat a_i\)，并再次检查结果。这样不会把``我已经改好了''一类依赖上文的句子当成第一次回答。

\section{离线算法：把纠正前移}

\subsection{构造只差一条纠正的两个输入}

对每条有效记录，构造
\begin{equation}
  x_i^{0}=x_i,\qquad
  x_i^{h}=x_i\oplus h_i.
\end{equation}
两个输入的任务、历史和工具状态完全相同，唯一差别是教师输入多了用户纠正。

在参考序列 \(\widehat a_i\) 的第 \(j\) 个位置，记录旧策略在两个输入上的分布：
\begin{align}
  p_{ij}^{0}(v)
    &=\pi_{\theta_{\mathrm{old}}}
      (v\mid x_i^{0},\widehat a_{i,<j}),\\
  q_{ij}(v)
    &=\sg\!\left[
      \pi_{\theta_{\mathrm{old}}}
      (v\mid x_i^{h},\widehat a_{i,<j})
      \right].
\end{align}
教师和学生使用同一个当前检查点；教师分支停止梯度。纠正带来的 token 方向为
\begin{equation}
  d_{ij}(v)
  =
  \clip\!\left(
    \log q_{ij}(v)-\log p_{ij}^{0}(v),-C,C
  \right).
  \label{eq:direction}
\end{equation}
共同的任务知识在相减时大体抵消，留下的是纠正改变了哪些 token 的概率。

\subsection{在没有纠正的输入上学习这个方向}

只在教师和学生 top-\(k\) token 的并集 \(S_{ij}\) 上计算，避免遍历整个词表。定义
\begin{align}
  w_{ij}(v)
    &=\operatorname{softmax}_{v\in S_{ij}}
      \log p_{ij}^{0}(v),\\
  A^{\mathrm{corr}}_{ij}(v)
    &=w_{ij}(v)d_{ij}(v),\\
  \rho_{ij}(v)
    &=\frac{\pi_{\theta}
       (v\mid x_i^{0},\widehat a_{i,<j})}
      {p_{ij}^{0}(v)}.
\end{align}
纠正前移损失为
\begin{equation}
  \mathcal L_{\mathrm{corr}}
  =
  \sum_{i,j}\sum_{v\in S_{ij}}
  \max\!\left[
    -A^{\mathrm{corr}}_{ij}(v)\rho_{ij}(v),
    -A^{\mathrm{corr}}_{ij}(v)
      \clip(\rho_{ij}(v),1-\epsilon,1+\epsilon)
  \right].
  \label{eq:corr-loss}
\end{equation}
这一步与 OpenClaw-RL 的 hint 蒸馏有直接关系
\cite{wang2026openclaw}。本文增加的是离线历史重建：从多轮日志找回
\(x_i^{0}\)，再把后续纠正配回这个状态。OPD 本身不是新贡献。

\subsection{同时学习真实任务结果}

纠正前移只告诉模型如何使用偏好，不能替代任务训练。对同一状态 \(x\)，保留或采样
一组真实动作：
\[
  \mathcal A_x=
  \{a^{-},\widehat a,a^{(1)},\ldots,a^{(K)}\}.
\]
执行每个动作得到
\[
  y_a=\FA.\operatorname{Execute}(x,a).
\]
任务分数 \(r_a^{\mathrm{task}}\in\{-1,+1\}\) 优先使用环境检查器。偏好分数
\(r_a^{\mathrm{pref}}\in[-1,1]\) 来自真实用户反应、可复现规则或 Judge。合成分数为
\begin{equation}
  r_a
  =r_a^{\mathrm{task}}
   +\lambda_{\mathrm{pref}}r_a^{\mathrm{pref}},
  \qquad 0\le\lambda_{\mathrm{pref}}<1.
  \label{eq:reward}
\end{equation}
因此，任务失败的动作不会靠写作风格超过任务成功的动作。报告时仍分别列出两项分数。

在同一状态的候选组内计算
\begin{equation}
  \bar r_x=\frac{1}{|\mathcal A_x|}
  \sum_{a\in\mathcal A_x}r_a,
  \qquad
  A_a=r_a-\bar r_x,
\end{equation}
再用现有的 clipped GRPO 得到
\(\mathcal L_{\mathrm{outcome}}\)。

\subsection{离线总目标}

\begin{equation}
  \mathcal L_{\mathrm{offline}}
  =
  \lambda_{\mathrm{corr}}\mathcal L_{\mathrm{corr}}
  +\lambda_{\mathrm{out}}\mathcal L_{\mathrm{outcome}}.
  \label{eq:offline}
\end{equation}
第一项把后续纠正学到早先状态；第二项保住答案和工具调用能力。两项都只更新
\(\lora\)。

\section{在线算法：继续用同一条更新}

离线训练得到 \(\pi_0\)。部署后，对每次交互保存
\[
  (x_t,a_t,y_t,f_t).
\]
若用户满意或任务成功，记录结果供 outcome 更新使用。若用户纠正：
\begin{enumerate}
  \item 冻结纠正发生前的 \(x_t\)；
  \item 从 \(f_t\) 提取 \(h_t\)；
  \item 构造 \(x_t^0=x_t\) 和 \(x_t^h=x_t\oplus h_t\)；
  \item 生成并检查参考动作 \(\widehat a_t\)；
  \item 用式~\eqref{eq:corr-loss} 在 \(x_t^0\) 上更新 LoRA；
  \item 同一批新轨迹用式~\eqref{eq:reward} 做 outcome 更新。
\end{enumerate}

在线阶段没有新的动作空间，也没有第二套优化器逻辑：
\begin{equation}
  \theta_{t+1}
  =
  \theta_t-\eta_t\nabla_\theta
  \left(
    \lambda_{\mathrm{corr}}\mathcal L_{\mathrm{corr}}^{(t)}
    +\lambda_{\mathrm{out}}\mathcal L_{\mathrm{outcome}}^{(t)}
  \right).
\end{equation}
区别只在数据：离线读取过去日志，在线读取刚发生的交互。

\section{完整工作流}

\subsection{离线}

\begin{enumerate}
  \item 按完整会话和时间把 \(\Dhist\) 切为 \(\Dtrain\) 与
    \(\Dholdout\)。先切分，再提取 \(P_u\) 或 hint。
  \item 从 \(\Dtrain\) 提取式~\eqref{eq:record} 的纠正记录。
  \item 过滤空反馈、损坏 hint 和未通过任务检查的记录。
  \item 对每条记录构造 \(x^0,x^h,\widehat a\)。
  \item 用式~\eqref{eq:corr-loss} 训练用户 LoRA。
  \item 在真实工具任务上采样动作并执行，用 GRPO 保住任务能力。
  \item 以固定训练预算得到 \(\pi_0\)。训练过程中不读
    \(\Dholdout\) 或 \(\Xholdout\)。
\end{enumerate}

\subsection{在线}

\begin{enumerate}
  \item 加载同一用户的 \(\pi_0\)。
  \item 先回答或执行任务，再读取真实 next state。
  \item 有纠正时，把纠正配回原状态；没有纠正时只保留结果信号。
  \item 累积到一个小批量后继续更新同一个 LoRA。
  \item 按预定间隔在开发任务上检查第一次回答和任务成功率。
  \item 最后只在一次冻结的评测中读取 \(\Dholdout,\Xholdout\)。
\end{enumerate}

\section{如何使用现有 OOCSO 数据}

\begin{table}[H]
  \centering
  \caption{第一轮可执行实验，不额外编造数据。}
  \begin{tabularx}{\textwidth}{@{}L{0.24\textwidth}L{0.31\textwidth}Y@{}}
    \toprule
    数据 & 进入算法的位置 & 说明 \\
    \midrule
    完整偏好会话 &
    \(\mathcal L_{\mathrm{corr}}\) &
    现有 dump 只用于单元测试；正式训练需要保存完整、可见的用户纠正 \\
    80 个离线工具任务 &
    \(\mathcal L_{\mathrm{outcome}}\) &
    沿用当前 \path{run.py} 的编号切分 \\
    40 个训练集中后段任务 &
    在线训练或开发实验 &
    在运行前固定用途，不能一边看结果一边改切分 \\
    30 个 dev 任务 &
    开发检查 &
    检查任务能力是否因个性化更新下降 \\
    30 个 eval 任务 &
    \(\Xholdout\) &
    最终只运行一次 \\
    \(\Xsim\) &
    暂不使用 &
    当前已有真实任务，先避免模拟问题引入新的偏好假设 \\
    \bottomrule
  \end{tabularx}
\end{table}

偏好实验和工具实验承担不同职责。偏好会话提供``用户希望怎样做''，工具任务提供
``任务有没有完成''。两者通过式~\eqref{eq:offline} 联合训练，但不能把工具任务自动标成某种
用户偏好。

\section{必须比较的版本}

\begin{table}[H]
  \centering
  \caption{最小对照。}
  \begin{tabularx}{\textwidth}{@{}L{0.31\textwidth}Y@{}}
    \toprule
    版本 & 回答的问题 \\
    \midrule
    \(\basepolicy\) &
    不做个性化时表现如何？ \\
    \(\basepolicy+P_u\) &
    每次把偏好放进提示词是否已经足够？ \\
    普通离线 GRPO &
    只用历史奖励、不给纠正前移时能学到什么？ \\
    普通离线 SFT/DPO &
    把纠正后回答当训练目标是否足够？ \\
    Online-only OpenClaw-RL/OPD &
    没有离线起点时需要多少次实时纠正？ \\
    本文：离线纠正前移 \(\rightarrow\) 在线继续更新 &
    后续纠正是否能提前作用于第一次回答？ \\
    \bottomrule
  \end{tabularx}
\end{table}

至少分开报告：
\begin{itemize}
  \item 第一次回答符合偏好的比例，以及连续三次符合所需的会话数；
  \item 工具任务成功率、平均步数和工具调用错误率；
  \item 各任务类别结果；
  \item 每位用户 LoRA 的大小、训练时间和在线更新次数。
\end{itemize}
不要把个性化和任务成功压成一个总分。

\section{创新边界}

以下部分已有工作已经覆盖：
\begin{itemize}
  \item 提示词优化：DSO、AGP 和 GEPA
    \cite{kiyohara2025prompt,wang2025agp,agrawal2025gepa}；
  \item 每位用户一个离线 LoRA：OPPU \cite{tan2024oppu}；
  \item 持续更新用户 LoRA：SPRInG \cite{kim2026spring}；
  \item 从 next state 提取 reward 和 hint，并做在线 OPD：
    OpenClaw-RL \cite{wang2026openclaw}；
  \item 离线到在线训练的一般过渡：已有半在线 DPO/RL 工作
    \cite{lanchantin2025bridging}。
\end{itemize}

因此，本文不能声称首次做提示词个性化、首次训练用户 LoRA，或首次使用 hint 蒸馏。
当前可检验的差异是：
\begin{quote}
从同一用户的历史中重建纠正前状态，把后来出现的纠正蒸馏回该状态，并在离线和部署后
使用同一条更新，同时用可执行任务结果约束动作。
\end{quote}
这仍是待验证的方法设计。实验成功和更完整的文献核对之前，不应写成已经成立的新算法。

\section{对现有代码的最小改动}

\begin{table}[H]
  \centering
  \caption{实现位置。本文只给设计，不在本次改动这些代码。}
  \begin{tabularx}{\textwidth}{@{}L{0.27\textwidth}Y@{}}
    \toprule
    文件 & 需要的后续改动 \\
    \midrule
    \path{pref/session.py} &
    保存纠正前状态、可见用户反馈、纠正后回答和是否被接受 \\
    \path{pref/pref_prm.py} &
    只从可见反馈提取 hint；继续过滤过短或损坏候选 \\
    \path{algo.py} &
    让 OPD 沿检查通过的参考序列计算，并使用 teacher/student top-\(k\) 并集 \\
    \path{pref/run_pref.py} &
    离线阶段加入纠正前移，不再只对带纠正的后续状态做 GRPO \\
    \path{agent.py} 与 \path{env/} &
    保持现有工具执行和确定性任务检查，不改工具协议 \\
    \bottomrule
  \end{tabularx}
\end{table}

第一项单元测试应直接使用现有苹果题记录：训练样本的 student 输入只能是第一次回答前的
\path{state_messages}，不能包含``别用格式''；teacher 输入才可以多出该纠正。若这一点没有
通过，算法仍会退回当前的``提醒后才改''。

{\footnotesize\raggedright
\begin{thebibliography}{9}
\setlength{\itemsep}{1pt}

\bibitem{wang2026openclaw}
Yinjie Wang, Xuyang Chen, Xiaolong Jin, Mengdi Wang, and Ling Yang.
\newblock OpenClaw-RL: Train Any Agent Simply by Talking.
\newblock \href{https://arxiv.org/abs/2603.10165}{arXiv:2603.10165}, 2026.

\bibitem{kiyohara2025prompt}
Haruka Kiyohara et al.
\newblock Prompt Optimization with Logged Bandit Data.
\newblock \href{https://arxiv.org/abs/2504.02646}{arXiv:2504.02646}, 2025.

\bibitem{wang2025agp}
Chen Wang et al.
\newblock Automating Personalization: Prompt Optimization for Recommendation
Reranking.
\newblock \href{https://arxiv.org/abs/2504.03965}{arXiv:2504.03965}, 2025.

\bibitem{agrawal2025gepa}
Lakshya A. Agrawal et al.
\newblock GEPA: Reflective Prompt Evolution Can Outperform Reinforcement
Learning.
\newblock \href{https://arxiv.org/abs/2507.19457}{arXiv:2507.19457}, 2025.

\bibitem{tan2024oppu}
Zhaoxuan Tan et al.
\newblock Democratizing Large Language Models via Personalized
Parameter-Efficient Fine-Tuning.
\newblock \href{https://aclanthology.org/2024.emnlp-main.372/}{EMNLP}, 2024.

\bibitem{kim2026spring}
Seoyeon Kim and Jaehyung Kim.
\newblock SPRInG: Continual LLM Personalization via Selective Parametric
Adaptation and Retrieval-Interpolated Generation.
\newblock \href{https://arxiv.org/abs/2601.09974}{arXiv:2601.09974}, 2026.

\bibitem{lanchantin2025bridging}
Jack Lanchantin et al.
\newblock Bridging Offline and Online Reinforcement Learning for LLMs.
\newblock \href{https://arxiv.org/abs/2506.21495}{arXiv:2506.21495}, 2025.

\end{thebibliography}
}

\end{document}
